\documentclass[aspectratio=169]{beamer}
\usetheme{Madrid}
\usecolortheme{default}
\usepackage{graphicx}
\usepackage{tabularx}
\usepackage{booktabs}
\usepackage{caption}

% Removes the navigation symbols at the bottom right
\setbeamertemplate{navigation symbols}{}

% Title Slide Customization
\title[Smart Street Light System]{IOT-BASED SMART STREET LIGHT SYSTEM\\ WITH CLOUD INTEGRATION}
\author[Your Name]{}
\institute[]{
    \vspace{0.5cm}
    \begin{columns}[t]
        \column{0.45\textwidth}
        \textbf{Under the Guidance of:}\\
        Dr. Faculty Name\\
        Assistant Professor\\
        Department of CSE
        
        \column{0.45\textwidth}
        \textbf{Submitted by:}\\
        Your Name\\
        Roll No.-22dcs026
    \end{columns}
    \vspace{1cm}
    
    \textbf{DEPARTMENT OF COMPUTER SCIENCE AND ENGINEERING}\\
    NATIONAL INSTITUTE OF TECHNOLOGY HAMIRPUR\\
    HIMACHAL PRADESH -177005
}
\date{Date: 01-04-2026}

\begin{document}

\begin{frame}
    \titlepage
\end{frame}

\begin{frame}{Table of Content}
    \begin{enumerate}
        \item[1.] Problem Statement
        \item[2.] Proposed Solution
        \item[3.] Methodology
        \begin{itemize}
            \item[3.1] Simulation Platforms (Wokwi \& TagoIO)
            \item[3.2] Sensors Used
            \item[3.3] Microcontroller
            \item[3.4] Circuit Connections
            \item[3.5] TagoIO Results
        \end{itemize}
        \item[4.] Limitations
        \item[5.] Future Scope
        \item[6.] Conclusion
    \end{enumerate}
\end{frame}

\begin{frame}{Problem Statement}
    \begin{columns}[T]
        \begin{column}{0.3\textwidth}
            \textbf{Energy Wastage:}
        \end{column}
        \begin{column}{0.65\textwidth}
            Traditional streetlights remain ON during daytime, wasting massive electricity globally.
        \end{column}
    \end{columns}
    \vspace{0.4cm}
    \begin{columns}[T]
        \begin{column}{0.3\textwidth}
            \textbf{Lack of Automation:}
        \end{column}
        \begin{column}{0.65\textwidth}
            Manual switching is prone to human errors and timing delays.
        \end{column}
    \end{columns}
    \vspace{0.4cm}
    \begin{columns}[T]
        \begin{column}{0.3\textwidth}
            \textbf{No Dimming Mechanism:}
        \end{column}
        \begin{column}{0.65\textwidth}
            Providing 100\% brightness even when roads are completely empty wastes power.
        \end{column}
    \end{columns}
    \vspace{0.4cm}
    \begin{columns}[T]
        \begin{column}{0.3\textwidth}
            \textbf{No Remote Monitoring:}
        \end{column}
        \begin{column}{0.65\textwidth}
            Faulty lights are only fixed after public complaints; no live dashboard.
        \end{column}
    \end{columns}
\end{frame}

\begin{frame}{Proposed Solution}
    A fully simulated, IoT-based Smart Street Light System that solves traditional limitations:
    \vspace{0.5cm}
    \begin{table}[]
        \begin{tabularx}{\textwidth}{@{}lX@{}}
            \toprule
            \textbf{Problem} & \textbf{Our Solution} \\
            \midrule
            \textbf{Energy Wastage} & Automated Day/Night switching using LDR sensor. \\ 
            \addlinespace
            \textbf{No Dimming} & PIR Motion Sensor dims lights to 20\% when empty, 100\% when active. \\ 
            \addlinespace
            \textbf{No Remote Access} & ESP32 WiFi pushes live telemetry data to the Cloud every few seconds. \\ 
            \addlinespace
            \textbf{Manual Operation} & Zero human intervention required. Fully autonomous. \\
            \bottomrule
        \end{tabularx}
    \end{table}
\end{frame}

\begin{frame}{Simulation Platforms Used}
    \textbf{Wokwi — Hardware Simulation}
    \begin{itemize}
        \item Free, browser-based IoT simulator at wokwi.com
        \item Simulates real ESP32 microcontroller with actual GPIO pins
        \item Supports LDR, PIR sensors, LEDs, Serial Monitor
        \item Arduino C++ code runs exactly as it would on physical hardware
    \end{itemize}
    \vspace{0.5cm}
    \textbf{TagoIO — Cloud Dashboard}
    \begin{itemize}
        \item Cloud IoT platform at tago.io
        \item Provides drag-and-drop dashboard widgets:
        \begin{itemize}
            \item LDR Chart
            \item Dimmer Widget
            \item Device Geolocation Map
            \item Analytics Dashboard
        \end{itemize}
    \end{itemize}
\end{frame}

\begin{frame}{Sensors Used}
    \textbf{1. LDR — Light Dependent Resistor (Photoresistor)}
    \begin{itemize}
        \item An analog sensor that changes resistance based on light intensity.
        \item Used to detect ambient environmental light (Day vs Night).
        \item Triggers the system to turn ON or OFF lights automatically based on natural sunlight.
    \end{itemize}
    \vspace{0.5cm}
    \textbf{2. PIR Sensor — Passive Infrared Sensor}
    \begin{itemize}
        \item A digital sensor that detects motion via infrared radiation changes.
        \item Used to identify presence of vehicles or pedestrians.
        \item Triggers the system to dynamically increase brightness to 100\% when motion is detected.
    \end{itemize}
\end{frame}

\begin{frame}{Microcontroller — ESP32}
    The ESP32 was selected as the central processing unit due to its high performance-to-cost ratio and native IoT capabilities.
    \vspace{0.3cm}
    \begin{itemize}
        \item \textbf{Dual-Core Processor:} Allows simultaneous sensor detection and Wi-Fi transmission.
        \item \textbf{Integrated Wi-Fi \& Bluetooth:} Enables direct MQTT/HTTP connection to TagoIO without external modules.
        \item \textbf{Rich I/O Peripheral Set:}
        \begin{itemize}
            \item ADC Pins: Used for reading variations in LDR resistance (Analog Input).
            \item Digital/PWM Pins: Used for controlling the Street Light LED brightness via PWM.
            \item Interrupt Pins: Used for real-time PIR motion detection.
        \end{itemize}
        \item \textbf{Low Power Consumption:} Essential for remote outdoor installations.
    \end{itemize}
\end{frame}

\begin{frame}{Circuit Architecture \& Connectivity}
    \textbf{LDR Sensor}
    \begin{itemize}
        \item VCC $\rightarrow$ 3.3V
        \item GND $\rightarrow$ GND
        \item AO $\rightarrow$ GPIO 34 (Analog Input)
    \end{itemize}
    \vspace{0.3cm}
    \textbf{PIR Sensor}
    \begin{itemize}
        \item VCC $\rightarrow$ 3.3V
        \item GND $\rightarrow$ GND
        \item OUT $\rightarrow$ GPIO 13 (Digital Input)
    \end{itemize}
    \vspace{0.3cm}
    \textbf{Street Light LED}
    \begin{itemize}
        \item Anode (+) $\rightarrow$ GPIO 25 (PWM Output)
        \item Cathode (-) $\rightarrow$ GND (via 220-ohm Resistor)
    \end{itemize}
\end{frame}

\begin{frame}{TagoIO Dashboard Results}
    \begin{figure}
        \centering
        \includegraphics[height=0.7\textheight,keepaspectratio]{dashboard_ss.png}
        \caption{Cloud Dashboard in TagoIO displaying live sensor states.}
    \end{figure}
\end{frame}

\begin{frame}{Limitations}
    \begin{columns}[T]
        \begin{column}{0.3\textwidth}
            \textbf{Simulation Constraints:}
        \end{column}
        \begin{column}{0.65\textwidth}
            Wokwi provides "perfect" data; real-world sensors face signal noise and physical degradation over time.
        \end{column}
    \end{columns}
    \vspace{0.5cm}
    \begin{columns}[T]
        \begin{column}{0.3\textwidth}
            \textbf{Network Dependency:}
        \end{column}
        \begin{column}{0.65\textwidth}
            System requires a constant Wi-Fi connection; lacks local fallback memory to operate completely offline without data loss.
        \end{column}
    \end{columns}
    \vspace{0.5cm}
    \begin{columns}[T]
        \begin{column}{0.3\textwidth}
            \textbf{Security:}
        \end{column}
        \begin{column}{0.65\textwidth}
            Data is sent over standard network protocols which might require enhanced MQTT encryption in enterprise solutions.
        \end{column}
    \end{columns}
\end{frame}

\begin{frame}{Future Scope}
    \begin{columns}[T]
        \begin{column}{0.3\textwidth}
            \textbf{Hardware Integration:}
        \end{column}
        \begin{column}{0.65\textwidth}
            Transition from Wokwi simulation to physical deployment using an ESP32 board, physical LDRs, and High-Power LEDs or Relays.
        \end{column}
    \end{columns}
    \vspace{0.5cm}
    \begin{columns}[T]
        \begin{column}{0.3\textwidth}
            \textbf{Edge AI / CV:}
        \end{column}
        \begin{column}{0.65\textwidth}
            Incorporate IP cameras to classify vehicles vs. pedestrians, overriding simple PIR triggers to predict trajectories.
        \end{column}
    \end{columns}
    \vspace{0.5cm}
    \begin{columns}[T]
        \begin{column}{0.3\textwidth}
            \textbf{Solar Panel Integration:}
        \end{column}
        \begin{column}{0.65\textwidth}
            Equip each streetlight pole with independent solar panels and battery storage to make the system completely self-sustaining.
        \end{column}
    \end{columns}
\end{frame}

\begin{frame}{Conclusion}
    \begin{columns}[T]
        \begin{column}{0.3\textwidth}
            \textbf{Integrated IoT Solution:}
        \end{column}
        \begin{column}{0.65\textwidth}
            Successfully demonstrated a real-time smart lighting system bridging ESP32 hardware logic with TagoIO cloud visualization.
        \end{column}
    \end{columns}
    \vspace{0.5cm}
    \begin{columns}[T]
        \begin{column}{0.3\textwidth}
            \textbf{Energy Efficiency:}
        \end{column}
        \begin{column}{0.65\textwidth}
            Moved beyond traditional 'always-on' lighting by implementing PIR-based dimming, actively drastic reducing power wastage.
        \end{column}
    \end{columns}
    \vspace{0.5cm}
    \begin{columns}[T]
        \begin{column}{0.3\textwidth}
            \textbf{Cost-Efficiency:}
        \end{column}
        \begin{column}{0.65\textwidth}
            The project proves that sophisticated urban automation can be achieved using low-cost, scalable, and off-the-shelf IoT components.
        \end{column}
    \end{columns}
\end{frame}

\end{document}
